\section{Charge Ionization Model (Shared by All Three Channels)}
\label{sec:ionization}

\subsection{Klein's formula: mean electron--hole pair creation energy}
\label{sec:klein}

Once a deposited energy $\omega\geq\Egap$ is fixed, the \textbf{mean} number
of electron--hole pairs created is set by the mean pair-creation energy
$\ehpair$ (historically called the ``Fano'' or ``ionization'' energy):
\begin{equation}
\langle n_e\rangle = \frac{\omega - \Egap}{\ehpair}\,.
\end{equation}
For real silicon, $\ehpair=3.8$~eV is \textbf{directly measured} (100~K,
PRD 102, 063026) — not derived from a formula. For a hypothetical material
with no measurement, \citet{Klein1968} proposed an empirical linear
relation between $\ehpair$ and the band gap, calibrated across several
elemental and compound semiconductors (Si, Ge, GaAs, CdS, \ldots):
\begin{equation}
\ehpair \approx a\,\Egap + b\,.
\end{equation}
CCDarkSens uses this class of relation with the coefficients
\begin{equation}
\boxed{\ehpair = 2.67\,\Egap + 0.87\ \text{eV}}
\end{equation}
which for \SrCdSb{} ($\Egap=0.34$~eV) gives
\begin{equation}
\ehpair = 2.67\times0.34 + 0.87 = 0.9078 + 0.87 = \mathbf{1.7778}\ \text{eV}\,.
\end{equation}
This value (\texttt{configs/darkphoton\_one\_point\_eu5in2sb6.json} and the
Sr2Cb2Sd scan configs) \textbf{supersedes} the earlier
$\ehpair=1.6$~eV placeholder used in the original HypMat unified-model
document. Note the repository also contains a second, coarser Klein-type
fit — the widely cited Alig--Bloom form $\ehpair = 2.8\Egap+0.5$~eV — used
only in exploratory heatmap plots
(\texttt{utils/plot\_band\_gap\_2d\_heatmap.py}); it is \textbf{not} the
parameterization used for the production \SrCdSb{} configs and is noted
here only to avoid confusion between the two curves if both appear in a
figure.

The actual $n_e$ for a given $\omega$ is \textbf{not sharp} — it follows a
Fano-broadened distribution $P(n_e\mid\omega)$, captured in a lookup table
(Section~\ref{sec:pne-table}), not by rounding $\langle n_e\rangle$.

\subsection{Ionization scenarios: the ``D-equal'' convention and its siblings}

Because there is \textbf{no known fundamental relation} $\ehpair(\Egap)$
valid for arbitrary low-gap materials (Klein's fit is itself only a
phenomenological average over a handful of real semiconductors),
CCDarkSens's band-gap phenomenology framework treats $(\Egap,\ehpair)$ as
two \textbf{independently specifiable} knobs and labels each combination
with a scenario ID:

\begin{center}
\begin{tabular}{@{}p{2cm}p{2.7cm}p{3.6cm}p{5.5cm}@{}}
\toprule
Scenario ID & $\Egap$ & $\ehpair$ & Interpretation \\
\midrule
\textbf{ref} & 1.2~eV & 3.8~eV & Reference Si, unscaled \texttt{p100K\_table.csv} \\
\textbf{B-thresh} & $\Egap$ (new) & 3.8~eV (fixed, Si-like) & Lower the threshold only, keep Si's pair scale \\
\textbf{D-equal} & $\Egap$ (new) & $=\Egap$ (new) & Aggressive equal-scale pheno; not real bulk Si \\
\textbf{A-ratio} & $\Egap$ (new) & $3.8\times\Egap/1.2$ & Preserve Si's $\ehpair/\Egap$ ratio \\
\textbf{Klein} & $\Egap$ (new) & $2.67\Egap+0.87$ & Section~\ref{sec:klein} relation; used for \SrCdSb{} \\
\textbf{C-grid} & user & user & Free 2D grid, no rule \\
\bottomrule
\end{tabular}
\end{center}

The \SrCdSb{} production configs use the \textbf{Klein} scenario
specifically (\texttt{ion\_scenario: "Klein"} in
\texttt{configs/Sr2Cb2Sd/write\_configs.py}), i.e.\ $\ehpair$ is
\emph{not} pinned equal to $\Egap$ (that would be D-equal) nor fixed at
Si's 3.8~eV (B-thresh) — it follows the Klein formula above.

\subsection{The $P(n_e\mid E)$ table and the anchored energy map}
\label{sec:pne-table}

The full ionization model is a table $P(n_e\mid E)$: CSV columns
$E_{r,\mathrm{eV}}, P_1, P_2, \ldots, P_k$, loaded by
\texttt{ChargeIonization::LoadCSV\_} and consumed by
\texttt{ChargeIonization::FoldToNe}, which convolves a $dR/dE$ histogram
against this table (weighted by exposure) to produce the observable $n_e$
spectrum.

For materials without a dedicated Monte Carlo $P(n_e\mid E)$ calculation,
the table is \textbf{rescaled from the Si reference table}
\texttt{data/p100K\_table.csv} ($\Egap=1.2$~eV, $\ehpair=3.8$~eV, 100~K) via
the anchored energy map:
\begin{equation}
E'(E) = \Egap^\mathrm{ref} + \bigl(E - \Egap^\mathrm{new}\bigr)\times
\frac{\ehpair^\mathrm{ref}}{\ehpair^\mathrm{new}} \qquad \text{for } E \geq \Egap^\mathrm{new}\,,
\end{equation}
\begin{equation}
P_\mathrm{new}(n_e\mid E) = \begin{cases}
0 & E < \Egap^\mathrm{new} \\
P_\mathrm{ref}\bigl(n_e \mid E'(E)\bigr) & \text{otherwise}
\end{cases}
\end{equation}
(linear interpolation on the reference table). \textbf{Physical
interpretation:} the map asks ``at what energy does bulk Si have the same
\emph{fractional} excitation above threshold as the new material at energy
$E$?'' — it anchors the ionization \emph{shape} (Fano broadening,
multiplicity structure) to Si's measured behavior while relocating the
\emph{threshold and scale} to the new material's $(\Egap,\ehpair)$.

For \SrCdSb{} ($\Egap=0.34$~eV, $\ehpair=1.7778$~eV via Klein):
\begin{equation}
E'(E) = 1.2 + (E-0.34)\times\frac{3.8}{1.7778} = 1.2 + (E-0.34)\times2.138\,,
\end{equation}
built by \texttt{utils/build\_p100K\_scaled.py --band-gap-eV 0.34
--eh-pair-eV 1.7778 --scenario Klein}.

\paragraph{What this does and does not claim:} the turn-on energy
(0.34~eV) and the mean-$n_e$ scale (via $\ehpair$) are physically motivated
(Klein's relation, however approximate); the detailed \emph{shape} of
$P(n_e\mid E)$ — Fano factor, sub-gap phonon tail — is inherited from Si
and is \textbf{not} a first-principles \SrCdSb{} calculation.
